\section{Introduction}
\vspace{-0.1cm}

\textit{Agents} \citep{wang2024survey}, language model-driven systems that operate through iterative cycles of reasoning, planning, and acting, adapting their behavior based on environmental or tool-generated feedback, have achieved strong performance in diverse applications, from code generation \citep{zhang2024codeagent, yang2024swe}, web browsing \citep{wei2025browsecomp, yao2022webshop}, medical decision-making \citep{heydari2025anatomy, mcduff2025towards, kim2024mdagents}, finance \citep{yu2025finmem}, sustainability \citep{zhang2025towards}, to scientific discovery \citep{gottweis2025towards, mitchener2025kosmos}. As tasks become more complex and require long-horizon environmental interaction, multi-agent systems (MAS) have gained attention as a way to support task decomposition, parallel exploration, and verification. At the same time, concurrent works question whether multi-agent coordination outperforms single-agent systems (SAS), leaving the conditions under which MAS provides genuine benefits remain underexplored~\citep{tran2025multiagent,guo2024large, cemri2025multi,gao2025single,openhands2025single,cognition2025multi}. Despite rapid adoption, there remains no principled quantitative framework for predicting when adding agents improves performance and when it instead introduces coordination costs that degrade it. This gap leaves practitioners relying on heuristics, hindering both the emergence of a science of agent systems and, critically for real-world deployment, the ability to determine when multi-agent coordination provides genuine value over simpler single-agent alternatives.

To determine when multi-agent coordination provides benefit, we first establish which task categories require agentic capabilities. 
A necessary prerequisite is distinguishing between \textit{agentic} and \textit{non-agentic} evaluation paradigms. Expanding from the Agentic Benchmark Checklist (ABC) introduced in \citep{zhu2025establishing}, we characterize \textit{agentic tasks} as those requiring: (i) sustained multi-step interactions with an external environment, (ii) iterative information gathering under partial observability, and (iii) adaptive strategy refinement based on environmental feedback.

These characteristics differentiate tasks like web browsing~\citep{wei2025browsecomp}, financial trading~\citep{yu2025finmem}, software engineering~\citep{jimenez2023swe}, and interactive planning~\citep{dagan2024plancraft} from traditional static benchmarks, tasks solvable through single-shot reasoning without environmental feedback, which lack external environments, are fully observed, or require identical solution strategies~\citep{liu2023agentbench,kapoor2024agents}. This distinction matters because, while recent agentic benchmarks have emerged (e.g., SWE-Bench~\citep{jimenez2023swe}, $\tau^2$-Bench~\citep{barres2025t2bench}, Terminal-Bench~\citep{merrill2026terminal}), \textit{multi-agent system evaluations} have been conducted predominantly on non-agentic tasks, potentially providing misleading guidance about when collaboration provides value. This distinction is practically consequential: while LLMs achieve high accuracy on isolated code generation tasks like HumanEval~\citep{chen2021humaneval}, real-world deployment requires agentic capabilities such as iterative debugging, repository navigation, and adaptive strategy refinement as exemplified by interactive coding assistants (e.g., Cursor, Copilot Workspace). Multi-agent systems that show monotonic improvement with team size on static benchmarks (reaching 89\% on HumanEval with five agents) exhibit fundamentally different scaling behavior when evaluated on tasks requiring sustained environmental interaction, where coordination overhead and error propagation dynamics dominate. 

At its core, this distinction reflects a trade-off between context integration and diversity~\citep{du2023improving, hong2024metagpt}. Single-agent systems maximize context integration by maintaining a unified memory stream in which all reasoning steps share full access to prior history, enabling effectively constant-time access to global context. In contrast, multi-agent systems impose intrinsic information fragmentation~\citep{tran2025multiagent}: while parallel agents enable diverse exploration, they incur an unavoidable \textit{coordination tax} in which the global context must be compressed into inter-agent messages. This lossy communication increases synchronization overhead and cognitive load~\citep{malone1994interdisciplinary}, fundamentally altering the scaling behavior of collaboration.

The underlying dynamics explain this discrepancy: on agentic tasks, coordination overhead scales with interaction depth, agents operate on progressively divergent world states, and errors cascade through execution chains rather than being corrected through voting. Recent work has identified cases where single strong models match or exceed multi-agent systems~\citep{gao2025single}, yet the evaluation literature provides limited guidance on \textit{what factors} determine collaborative success, whether semantic diversity predicts team performance, how architectural choices shape coordination costs, or whether agents can detect and correct failures in extended interactions.

The problem is further compounded by rapid progress in frontier model capabilities. As base LLMs gain extended context windows, sophisticated tool use, and improved self-reflection, the unique value proposition of multi-agent collaboration becomes unclear. The answer likely depends on task characteristics and architectural choices that remain to be systematically quantified.

Two key challenges hinder progress toward principled multi-agent design. \textbf{First}, existing MAS evaluations compare architectures using different prompts, tools, or computational budgets, conflating architectural effects with implementation choices and precluding clean causal attribution. \textbf{Second}, evaluations focus exclusively on final accuracy metrics without examining process dynamics such as coordination overhead, error propagation, and information flow that determine whether collaboration succeeds or fails. We know from human team performance~\citep{mcgrath1964,lencioni2002} that team effectiveness depends on composition, coordination mechanisms, and member differentiation. Yet we lack comparable empirical understanding of how these principles translate to artificial agents, leaving practitioners without quantitative guidance for architecture selection.

To address these challenges, we present a controlled evaluation establishing the principles for agent coordination. Our experimental design isolates architectural effects by controlling for implementation confounds which maintains identical task prompts, tools, and computational budgets across all configurations, while systematically varying only coordination structure and model capability. \textcolor{black}{We evaluate five canonical architectures: Single Agent System (SAS) and four Multi-Agent variants (Independent, Centralized, Decentralized, Hybrid) instantiated across three major LLM families (OpenAI, Google, Anthropic) sampling models at varying capability tiers as quantified by an aggregate Intelligence Index (see Appendix \ref{sec:appendix-intelligence-index}), on six agentic benchmarks: (1) web browsing (BrowseComp-Plus \citep{chen2025browsecomp}), (2) financial analysis (Finance-Agent \citep{bigeard2025finance}), (3) game planning (PlanCraft \citep{dagan2024plancraft}), (4) realistic workplace tasks (Workbench \citep{styles2024workbench}), (5) software engineering (SWE-bench Verified \citep{jimenez2023swe}), and (6) terminal tasks (Terminal-Bench \citep{merrill2026terminal}). Across $N{=}260$ controlled configurations with matched compute, we derive a scaling principle across tested domains quantifying how performance emerges from empirically measured coordination properties.}

In contrast to prior claims that ``\textit{more agents is all you need}'', our evaluation reveals that the effectiveness of multi-agent systems is governed by quantifiable trade-offs between architectural properties and task characteristics. We establish a predictive framework using a mixed-effects regression model with empirical coordination metrics: efficiency (success/overhead ratio), error amplification factors, message density and redundancy as predictors, \textcolor{black}{achieving cross-validated $R^2{=}0.373$ across all six benchmarks ($R^2{=}0.413$ with a task-grounded capability metric), without dataset-specific parameters.} Critically, this framework generalizes beyond the fitted configurations in a restricted sense, where it predicts the best-performing architecture for 87\% of held-out task configurations, indicating relative architecture selection is more stable than absolute cross-domain performance prediction.

\textcolor{black}{Our analysis identifies three scaling patterns.} First, a \textit{tool-coordination trade-off} ($\beta{=}{-}0.096$, $p{=}0.002$): tool-heavy tasks (e.g., 16-tool business workflows) suffer from multi-agent coordination overhead, with efficiency penalties compounding as environmental complexity increases. Second, a \textit{capability ceiling} ($\beta{=}{-}0.236$, $p{=}0.004$): tasks where single-agent performance already exceeds 45\% accuracy experience negative returns from additional agents, as coordination costs exceed diminishing improvement potential. 
Third, we observe \textit{architecture-dependent error amplification}. Independent systems amplify trace-level errors $17.2\times$ through \textit{unchecked error propagation}, where individual mistakes cascade to the final output. Centralized coordination, however, contains this to $4.4\times$ by enforcing \textit{validation bottlenecks} that intercept errors before aggregation. Performance spans ${+}80.8\%$ relative improvement (structured financial reasoning under centralized coordination) to ${-}70.0\%$ degradation (sequential planning under independent coordination), demonstrating that architecture-task alignment, not number of agents, determines collaborative success. Optimal architectures vary systematically: decentralized coordination benefits tasks requiring parallel exploration of high-entropy search spaces (dynamic web navigation: ${+}9.2\%$), while all multi-agent variants universally degrade performance on tasks requiring sequential constraint satisfaction (planning: ${-}39\%$ to ${-}70\%$), where coordination overhead fragments reasoning capacity under fixed computational budgets. We translate these findings into quantitative architecture selection 
rules (Section~\ref{subsec:scaling}) achieving 87\% prediction accuracy on held-out configurations. The underlying mechanisms driving these patterns are interpretable: the tool-coordination trade-off arises because multi-agent systems fragment the per-agent token budget, leaving insufficient capacity for complex tool orchestration; the capability 
ceiling reflects that coordination overhead becomes a net cost when 
baseline performance is already high; and architecture-dependent error amplification stems from the presence or absence of validation bottlenecks that catch errors before propagation. These mechanistic insights enable 
practitioners to move from architectural heuristics to principled, measurement-driven deployment decisions.

Our primary contributions are:

\begin{itemize}

\item \textbf{Controlled evaluation of agent systems:} We establish a framework for comparing agent architectures, controlling for implementation confounds to isolate the effects of coordination structure. \textcolor{black}{Our framework spans 260 configurations across three LLM families and six diverse benchmarks, enabling controlled attribution of performance differences to architectural choices rather than stochastic variations.}

\item \textbf{Intelligence-Coordination alignment:} We characterize the non-linear relationship between foundational model capabilities and agentic performance. We demonstrate that while higher capability (Intelligence Index) yields consistent linear returns, these gains are not automatic; they strictly depend on architectural alignment. Without correct coordination structures, foundational improvements are often negated by coordination overhead.

\item \textbf{Quantitative scaling principles and architecture alignment:} We derive a regression model (\textcolor{black}{$R^2{=}0.373$ across all six benchmarks; $R^2{=}0.413$ with a task-grounded capability metric}) using empirical coordination metrics, efficiency ($E_c$), trace-level error amplification ($A_e^{\text{trace}}$), and redundancy ($\rho$) to quantify how performance emerges from the interplay of reasoning capability and task properties. This framework identifies fundamental limits on coordination, specifically a \textit{tool-coordination trade-off} ($\beta{=}{-}0.096$) where tool-heavy workflows suffer from coordination tax, and safety bounds where centralized verification reduces trace-level error amplification from $17.2\times$ to $4.4\times$. Using these mechanisms, we demonstrate that architecture selection is governed by measurable task features (e.g., decomposability) rather than simple agent scaling, achieving 87\% accuracy in predicting optimal architectures on held-out tasks.
\end{itemize}

